Für das Grafikdesign werden verschiedene Bilder im Pixelartdesign benötigt, welche dann im Spiel zu sehen sind.
Der Spieler soll diese Hintergründe per Point-and-Click ändern können, sodass er die Spielfigur im Schulhaus bewegen kann. \\ 
Dafür ist es wichtig, dass die verschiedenen Hintergründe so miteinander verknüpft werden, dass der Aufbau der Szenen der realen Struktur unserer Schule ähnelt.
Dies soll so funktionieren, dass es einen Szenenstapel gibt, in dem einzelne Szenen mithilfe von Objekten, auf die man klicken kann, ersetzt werden können. 
Für jede Szene wird eine eigene Klasse erstellt, in der die jeweils benötigten Objekte eingefügt werden. \\
Eine weitere Aufgabe der Gruppe bestand darin, dafür zu sorgen, dass die Spielfigur in den Szenen auftaucht, wobei auf eine vielfältige Variabilität der Position
der Figur zu achten ist. \\ 
Die Spielfigur, welche als Objekt vorliegt, wird in jeder Szene individuell hinzugefügt.
Zur Umsetzung von verschiedenen Positionen der Spielfigur könnte man eine Methode entwerfen, welche bei jedem Szenenwechsel eine zufällige x-Koordinate für die 
Spielfigur erstellt. Die y-Koordinate sollte fest definiert sein, da unsere Spielfigur nicht auf dem Bildschirm \glqq schweben\grqq\ soll.
